\documentclass[11pt,a4paper]{article}

%%%%%%%%%%%%%%%%%%%%%%%%%%%%%%%%%%%%%%%%%%%%%%%%%%%%%%%%%%%%%%%%%%%%%%%%
% Résumé substantiel en français (EDITE / Sorbonne Université)
% Document isolé, réutilise bibliography.bib du manuscrit
% Structure : par contribution (option A)
%%%%%%%%%%%%%%%%%%%%%%%%%%%%%%%%%%%%%%%%%%%%%%%%%%%%%%%%%%%%%%%%%%%%%%%%

\usepackage[utf8]{inputenc}
\usepackage[T1]{fontenc}
\usepackage[french]{babel}
\usepackage{lmodern}
\usepackage{microtype}
\usepackage[a4paper,margin=2.5cm]{geometry}
\usepackage{setspace}
\onehalfspacing

\usepackage{xspace}
\usepackage{graphicx}
\usepackage{booktabs}
\usepackage{multirow}
\usepackage{array}
\usepackage{amsmath,amssymb}
\usepackage{siunitx}
\sisetup{output-decimal-marker={,},group-separator={\,}}

\usepackage[round]{natbib}
\usepackage{xcolor}
\usepackage[%
  unicode,
  colorlinks=true,
  citecolor=blue,
  linkcolor=blue,
  urlcolor=blue,
]{hyperref}

% Métadonnées
\title{\textbf{Résumé substantiel}\\[0.4em]
\large Génération et évaluation de données médicales synthétiques\\
pour le traitement automatique des langues cliniques}
\author{Simon Meoni\\
\small Sorbonne Université -- École Doctorale Informatique, Télécommunications et Électronique (EDITE)}
\date{\today}

\begin{document}
\maketitle

\begin{abstract}
\noindent
Ce document constitue le résumé substantiel en français du manuscrit de thèse rédigé en anglais,
conformément aux exigences de l'EDITE de Paris pour les doctorants de Sorbonne Université.
La thèse porte sur la génération de données cliniques synthétiques préservant la vie privée
et sur leur évaluation pour des tâches de traitement automatique des langues médicales.
\end{abstract}

\tableofcontents
\vspace{1em}
\hrule
\vspace{1em}

%%%%%%%%%%%%%%%%%%%%%%%%%%%%%%%%%%%%%%%%%%%%%%%%%%%%%%%%%%%%%%%%%%%%%%%%
\section{Introduction et problématique}
\label{sec:introduction}
%%%%%%%%%%%%%%%%%%%%%%%%%%%%%%%%%%%%%%%%%%%%%%%%%%%%%%%%%%%%%%%%%%%%%%%%

\subsection{Contexte et verrou de la vie privée}
% Contexte TAL clinique, RGPD/HIPAA, rareté des corpus annotés

\subsection{Question de recherche}
% Comment générer des documents cliniques synthétiques utiles pour des
% tâches en aval tout en maîtrisant le risque de réidentification ?

\subsection{Contributions et plan du résumé}
% 1) Génération fac-similé guidée par mots-clés UMLS (chap. 3)
% 2) Alignement itératif par DPO avec scoreur contrastif privé (chap. 4)
% 3) Transfert au français clinique AP-HP, CIM-10 (chap. 5)
% 4) Annotation faible multilingue par distillation LLM (chap. 6)
% 5) Analyse vie privée : attaques par liaison et de proximité (chap. 7)
% 6) Stratégies d'atténuation et arbitrage utilité/confidentialité

%%%%%%%%%%%%%%%%%%%%%%%%%%%%%%%%%%%%%%%%%%%%%%%%%%%%%%%%%%%%%%%%%%%%%%%%
\section{État de l'art}
\label{sec:related}
%%%%%%%%%%%%%%%%%%%%%%%%%%%%%%%%%%%%%%%%%%%%%%%%%%%%%%%%%%%%%%%%%%%%%%%%

\subsection{Traitement automatique des langues cliniques}
% MIMIC-III, corpus cliniques, tâches (NER, classification CIM)

\subsection{Génération de données cliniques synthétiques}
% Approches par LLM, prompt-based, distillation, KnowledgeSG

\subsection{Évaluation de la confidentialité}
% Dé-identification, confidentialité différentielle, attaques d'inférence

\subsection{Positionnement : fac-similé vs fictionnel}
% Frontière de confidentialité, trade-off utilité/confidentialité

%%%%%%%%%%%%%%%%%%%%%%%%%%%%%%%%%%%%%%%%%%%%%%%%%%%%%%%%%%%%%%%%%%%%%%%%
\section{Génération fac-similé guidée par mots-clés médicaux}
\label{sec:facsimile}
%%%%%%%%%%%%%%%%%%%%%%%%%%%%%%%%%%%%%%%%%%%%%%%%%%%%%%%%%%%%%%%%%%%%%%%%
% Source : sources/chapters/synthetic_generation_cl4health_2025.tex

\subsection{Principe de la génération fac-similé}
% Définition, intuition, frontière de confidentialité,
% extraction UMLS comme seul signal traversant la frontière

\subsection{Protocole expérimental sur MIMIC-III}
% Extraction keywords, prompting, choix du générateur,
% construction d'AlpaCare

\subsection{Résultats : classification CIM-9}
% Performances comparables ou supérieures aux données réelles,
% effet de l'oversampling sur espaces de labels complexes

%%%%%%%%%%%%%%%%%%%%%%%%%%%%%%%%%%%%%%%%%%%%%%%%%%%%%%%%%%%%%%%%%%%%%%%%
\section{Alignement itératif par DPO}
\label{sec:dpo}
%%%%%%%%%%%%%%%%%%%%%%%%%%%%%%%%%%%%%%%%%%%%%%%%%%%%%%%%%%%%%%%%%%%%%%%%
% Source : sources/chapters/synthetic_generation_improvements.tex

\subsection{Motivation : écart de qualité résiduel}
% Limites de la génération one-shot, nécessité d'un alignement

\subsection{Scoreur contrastif privé et boucle DPO}
% Architecture du scoreur, préférences, boucle itérative,
% rôle de la frontière de confidentialité

\subsection{Résultats : gains d'utilité}
% Courbes d'alignement, gains CIM-9, robustesse

%%%%%%%%%%%%%%%%%%%%%%%%%%%%%%%%%%%%%%%%%%%%%%%%%%%%%%%%%%%%%%%%%%%%%%%%
\section{Transfert au français clinique : AP-HP et CIM-10}
\label{sec:french}
%%%%%%%%%%%%%%%%%%%%%%%%%%%%%%%%%%%%%%%%%%%%%%%%%%%%%%%%%%%%%%%%%%%%%%%%
% Source : sources/chapters/french_clinical_application.tex

\subsection{Contexte hospitalier et données AP-HP}
% Spécificités du français clinique, contraintes institutionnelles

\subsection{Adaptation du pipeline de génération}
% Transfert du protocole, adaptation CIM-10

\subsection{Résultats et preuve de généralisation}
% Gains d'utilité, confirmation de l'indépendance linguistique

%%%%%%%%%%%%%%%%%%%%%%%%%%%%%%%%%%%%%%%%%%%%%%%%%%%%%%%%%%%%%%%%%%%%%%%%
\section{Annotation faible multilingue par distillation LLM}
\label{sec:weak}
%%%%%%%%%%%%%%%%%%%%%%%%%%%%%%%%%%%%%%%%%%%%%%%%%%%%%%%%%%%%%%%%%%%%%%%%
% Source : sources/chapters/weak_annotations.tex

\subsection{Verrou : NER clinique en contexte faiblement doté}
% E3C, cinq langues européennes, rareté des annotations gold

\subsection{Supervision hybride : dictionnaires et LLM}
% Distillation d'un LLM vers des modèles encodeurs,
% complémentarité dictionnaire + LLM

\subsection{Résultats multilingues}
% Gains pour les langues faiblement dotées, comparaison des stratégies

%%%%%%%%%%%%%%%%%%%%%%%%%%%%%%%%%%%%%%%%%%%%%%%%%%%%%%%%%%%%%%%%%%%%%%%%
\section{Analyse de la vie privée et atténuation}
\label{sec:privacy}
%%%%%%%%%%%%%%%%%%%%%%%%%%%%%%%%%%%%%%%%%%%%%%%%%%%%%%%%%%%%%%%%%%%%%%%%
% Source : sources/chapters/discussion.tex (section vie privée)

\subsection{Modèle de menace et protocoles d'attaque}
% Attaques par liaison TF-IDF, attaques de proximité, discrimination

\subsection{Résultats : risques de réidentification}
% Linkage 45--76\%, discrimination 89--96\%,
% effet amplificateur du DPO sur le risque

\subsection{Stratégies d'atténuation}
% Perturbation des mots-clés, effet négligeable sur l'utilité conversationnelle

\subsection{Comparaison à KnowledgeSG}
% Complétion d'identifiants vs réidentification,
% complémentarité des deux approches

%%%%%%%%%%%%%%%%%%%%%%%%%%%%%%%%%%%%%%%%%%%%%%%%%%%%%%%%%%%%%%%%%%%%%%%%
\section{Discussion et conclusion}
\label{sec:conclusion}
%%%%%%%%%%%%%%%%%%%%%%%%%%%%%%%%%%%%%%%%%%%%%%%%%%%%%%%%%%%%%%%%%%%%%%%%

\subsection{Arbitrage utilité/confidentialité}
% Tensions fondamentales, pas de risque nul, honnêteté scientifique

\subsection{Limites}
% Annotations silver, dépendance au contexte institutionnel, généralisation

\subsection{Perspectives}
% Frameworks de remédiation, métriques composites, alignement réglementaire,
% applications hors médical (juridique, financier)

%%%%%%%%%%%%%%%%%%%%%%%%%%%%%%%%%%%%%%%%%%%%%%%%%%%%%%%%%%%%%%%%%%%%%%%%
% Bibliographie : réutilise bibliography.bib du manuscrit principal
%%%%%%%%%%%%%%%%%%%%%%%%%%%%%%%%%%%%%%%%%%%%%%%%%%%%%%%%%%%%%%%%%%%%%%%%
\bibliographystyle{styles/acl_natbib}
\bibliography{bibliography}

\end{document}
